\subsection{I : objets, variance et domaines}
\label{note:base1-functorial}
Le spectre est contravariant (1.2.3) : à une composition
$A\to B\to C$ correspond $\operatorname{Spec}(C)\to
\operatorname{Spec}(B)\to\operatorname{Spec}(A)$.
Le point générique de 2.1.5 est celui de $Y$.
En 2.5.5, l'objet est un $S$-préschéma ; en 3.3.8, la flèche considérée
est $g$. Le but du changement de base de 3.3.12 est $X'$.
En 3.3.15, le morphisme vers l'espace affine a le sens
$X\to\operatorname{Spec}(\mathbf Z[T])$.
En 3.7.2, l'adhérence contient $X'$.

Pour une immersion en 4.2.2, le morphisme de faisceaux d'anneaux est
la surjection $\psi^*\mathcal O_X\to\mathcal O_Y$ ; l'application sur
les sections de 4.2.3 provient de $\theta$.
L'extension de l'idéal $K\subset B$ à $A$ en 4.4.5 utilise la structure
de $B$-algèbre de $A$. L'argument sur les corps résiduels de 4.5.5
se fait pour un point $z$ et son image $z'$.
La restriction de 5.1.9.2 est $\mathcal F|V$.
Le complément employé dans l'ouvert $U$ en 6.1.8 doit être intersecté
avec $U$. Les voisinages de 6.6.2 et 6.6.3 sont ceux des points images
$f(x)$ et $y$ respectivement. En 9.5.11, le schéma considéré après
changement de base est $Y'$.

\subsection{I, 3--4 : produits, sommes et quotients}
\label{note:base1-categorical}
Les propriétés universelles des produits de 3.3 s'interprètent dans
une catégorie où ces produits existent. L'assertion de 3.3.10 sur
les sommes exige cependant une propriété supplémentaire de la
catégorie. Dans les ensembles pointés, si $A$ a deux éléments et
$*=\{*\}$, alors $A\times(*\vee*)=A$ a deux éléments, tandis que
$(A\times*)\vee(A\times*)$ en a trois. Le produit ne distribue
donc pas toujours les sommes. Cette assertion appartient aux
exceptions énumérées en 3.3.1.

En 3.4.9, le morphisme d'un produit tensoriel vers un corps résiduel
n'est pas nécessairement injectif. La multiplication
$\mathbf C\otimes_{\mathbf R}\mathbf C\to\mathbf C$ fournit
un exemple : $i\otimes1-1\otimes i$ est non nul et a une image nulle.
En 4.3.1, pour des idéaux $\mathfrak b\subset B$ et
$\mathfrak c\subset C$, le noyau du morphisme vers
$(B/\mathfrak b)\otimes_A(C/\mathfrak c)$ est
\[
 \mathfrak b(B\otimes_A C)+\mathfrak c(B\otimes_A C).
\]
Les images de $\mathfrak b$ et $\mathfrak c$ engendrent ces idéaux ;
leurs simples images additives ne constituent pas nécessairement
des idéaux du produit tensoriel.

\subsection{I, 5.3 et 6.1 : diagonale et non-vacuité}
\label{note:base1-diagonal}
Dans 5.3.5, le morphisme structural de $Y$ est $g:Y\to S$.
Pour un sous-foncteur du foncteur final, chaque ensemble de points
contient \emph{au plus} un élément (5.3.8) : l'injection dans un
singleton permet aussi l'ensemble vide.

La preuve de 5.3.9 est fondée sur l'erratum des auteurs publié dans
III, seconde partie, p.~220, Errata III, no~10.
Pour la vérification locale, on choisit un ouvert affine $V\subset S$
et un ouvert affine $U\subset f^{-1}(V)$. Alors
$U\times_S U=U\times_V U$ est affine, et la diagonale y correspond
à la surjection $A\otimes_B A\to A$, $a\otimes a'\mapsto aa'$.
Elle est donc une immersion fermée sur cet ouvert. Ces ouverts
couvrent l'image de la diagonale ; le critère local de I, 4.2.4(a)
donne l'immersion. Le choix explicite de $V$ précise le cadre affine
de l'argument des auteurs.

Dans 6.1.12, la caractérisation de l'intégrité exige un espace
connexe \emph{non vide} en plus des conditions sur ses anneaux locaux.
L'espace vide satisferait sinon vacuement les conditions locales.

\subsection{I, 7 : applications rationnelles et rang générique}
\label{note:base1-rational}
Deux représentants d'une même application rationnelle coïncident
sur un ouvert dense commun. En 7.1.2, la comparaison des trois
représentants utilise donc un ouvert dense de $U\cap V\cap W$,
sans affirmer leur égalité sur toute cette intersection.

Pour I, 7.3.8, la lecture est celle de l'erratum des auteurs publié
dans II, pp.~221--222. Pour $X,Y$ intègres et $f:X\to Y$ dominant,
on dispose du morphisme canonique
$\tau:f^*\mathscr R(Y)\to\mathscr R(X)$, qui est un isomorphisme
au point générique de $X$. Il ne s'agit pas d'un isomorphisme global
en général : pour $X=\mathbb A^1_k$ et $Y=\operatorname{Spec}(k)$,
le morphisme est l'inclusion $\mathcal O_X\to k(t)$.

Le sens réciproque de 7.4.3 porte sur un sous-module
\emph{quasi-cohérent non nul} de $\mathscr R(X)$.
Il est sans torsion et son germe générique est un sous-espace non nul
du corps des fonctions, donc de dimension un sur ce corps.
Le sous-module nul aurait rang zéro ; la quasi-cohérence ne résulte
pas de la seule inclusion dans un faisceau.
En 7.4.4, les germes génériques sont des espaces vectoriels sur
le corps des fonctions : le produit tensoriel de deux vecteurs
non nuls y est non nul. En 7.4.7, pour un espace réductible, la
condition sur le support exclut \emph{chaque} composante irréductible,
ce qui exprime l'annulation à chacun des points génériques.

\subsection{I, 9.4--9.5 : recollement et fermeture schématique}
\label{note:base1-transfinite}
Pour le recouvrement bien ordonné de 9.4.6, soit $\kappa$ son type
d'ordre, avec les ouverts $V_\mu$ pour $\mu<\kappa$.
On pose $U_\lambda=U\cup\bigcup_{\mu<\lambda}V_\mu$ pour
$0\leq\lambda\leq\kappa$. L'étape initiale est le module donné
$\mathcal G$ sur $U_0=U$. Une étape successeur étend le module
à l'ouvert suivant par l'argument local de la proposition.
À un ordinal limite, les modules compatibles se recollent sur
la réunion ; la quasi-cohérence et le type fini sont des propriétés
locales. L'étape $\kappa$ est indispensable, y compris lorsque le
recouvrement a un dernier indice, car $U_\kappa=X$.
On obtient ainsi le module voulu sur tout $X$ sous les hypothèses
de la proposition.

Dans 9.5.1, la fermeture schématique est le plus petit
\emph{sous-schéma fermé} ayant la propriété considérée.
La condition de fermeture est partie intégrante de cette
caractérisation par un idéal quasi-cohérent.

\subsection{I, 10 : schémas formels, idéaux et adjonction}
\label{note:base1-formal}
Pour un idéal ouvert $J$ de $A$, le spectre formel de $A/J$ s'identifie
au fermé $V(J)$ de $\operatorname{Spf}(A)$ en 10.3.2.
Ce fermé n'est pas nécessairement l'espace entier : pour $k[t]$ discret
et $J=(t)$, il est le point $(t)$. Pour un idéal de définition,
le fermé contient tous les premiers ouverts et est l'espace entier.

La section $h$ de 10.9.4 appartient à $\Gamma(X,\mathcal O_X)$.
Les rangs des modules libres dans 10.10.5 sont $p,q$.
La correspondance de 10.14.2 concerne les idéaux cohérents : le
quotient par un idéal est une algèbre et définit un sous-schéma fermé.
En 10.14.4, hors de l'image fermée, l'idéal est égal à $\mathcal O_X$ ;
c'est le quotient par cet idéal qui a le support fermé annoncé.
La factorisation est $f=j\circ g$, dans l'ordre des domaines.
Enfin, le morphisme $u$ de 10.14.6 est l'adjoint de $f_*(\rho)$,
où $\rho:\mathcal F\to p_*p^*\mathcal F$ est l'unité d'adjonction.
Les deux homomorphismes appartiennent à des groupes Hom différents,
identifiés par l'adjonction ; ils ne sont pas littéralement la même flèche.

\subsection{I, 10.9.9 : l’objet formel}
\label{note:base1-formalname}
Le complété $X_{/X'}$ est un préschéma formel.
Cette qualification indique la catégorie dans laquelle est pris
le diagramme de 10.9.9 et la structure topologique de ses anneaux.

\subsection{I : repères des critères locaux et des compositions}
\label{note:base1-refs}
La définition invoquée dans 10.5.1 est 10.3.7.
Le résultat sur la composition des immersions, utilisé dans 4.5.5
et 5.3.13, est 4.2.5. Le renvoi de 5.1.4 est 2.1.8 ; celui de
5.4.8 est 5.4.4. En particulier le critère local d'immersion
de 4.2.4 et sa stabilité par composition en 4.2.5 ont des rôles distincts.
\par
